\section{Risultati}
%  \begin{frame}
% \frametitle{Distribuzione di flusso neutronico}
% \begin{figure}[!htb]
% \centering
% \subfigure[Sezione con barre inserite]
% {\includegraphics[height=0.28\textheight]{../../thesis/figure/result_tred/morenew/pw-clip-insert1.png}}
% \hspace{1.2cm}
% \subfigure[Nocciolo con barre inserite]
% {\includegraphics[height=0.28\textheight]{../../thesis/figure/result_tred/morenew/pw-trashold-insert1.png}}\\
% \subfigure[Sezione con barre estratte]
% {\includegraphics[height=0.28\textheight]{../../thesis/figure/result_tred/morenew/pw-clip-exstract1.png}}
% \hspace{1.2cm}
% \subfigure[Nocciolo con barre estratte]
% {\includegraphics[height=0.28\textheight]{../../thesis/figure/result_tred/morenew/pw-trashold-exstract1.png}}
% \end{figure}
% \end{frame}
%  \begin{frame}
% \frametitle{Distribuzione di flusso neutronico}
% \begin{figure}[!htb]
% \centering
% \subfigure[Sezione con barre inserite]
% {\includegraphics[height=0.28\textheight]{../../thesis/figure/result_tred/morenew/pw-clip-insert1.png}}
% \hspace{1.2cm}
% \subfigure[Nocciolo con barre inserite]
% {\includegraphics[height=0.28\textheight]{../../thesis/figure/result_tred/morenew/pw-trashold-insert1.png}}\\
% \subfigure[Sezione con barre estratte]
% {\includegraphics[height=0.28\textheight]{../../thesis/figure/result_tred/morenew/pw-clip-exstract1.png}}
% \hspace{1.2cm}
% \subfigure[Nocciolo con barre estratte]
% {\includegraphics[height=0.28\textheight]{../../thesis/figure/result_tred/morenew/pw-trashold-exstract1.png}}
% \end{figure}
% \end{frame}
% 
\begin{frame}
\frametitle{Validazione Neutronica}
\vspace{0.4cm}
\begin{center}

  Confronto tra i valori della costante moltiplicativa effettiva $k_{eff}$ 
  tra il codice deterministico \textbf{\textcolor{blue}{Donjon}} e il codice probabilistico
  \textbf{MCNPX} effettuato in un altro studio
  
  
 \begin{center}
 \begin{tabular}{l|cc|cc}
\hline
Tempo [Giorni] & $k_{eff}$ Donjon [-] & & $k_{eff}$ MCNPX [-] & \\
$t_{burnup}$ & Barre Estratte & Inserite & Barre Estratte & Inserite \\
\hline
0   & 1.1177  & 1.0794  & 1.0804 & 1.0525  \\
365 & 1.0887  & 1.0520  & 1.0510 & 1.0229   \\
730 & 1.0624  & 1.0272  & 1.0247 & 0.9964   \\
1095 & 1.0384  & 1.0001 & 0.9988  & 0.9703 \\
\hline
\end{tabular}
  \end{center}
  \begin{figure}
   \includegraphics[scale=0.15]{../barre-di-controllo}
  \end{figure}
 inserite\quad \quad  \quad estratte
\end{center}  
\end{frame}
% %    %%%%%%%%%%%%%%%%%%%%%%%%%%%%%%%%%%%%%%%%%%%%%%%%%%%%%%%%%%%%%%
% % \begin{frame}
% % \frametitle{Validazione Neutronica}
% % Confronto tra i valori della costante moltiplicativa effettiva $k_{eff}$ 
% %   tra il caso deterministico di \textcolor{blue}{Donjon} e il caso probabilistico
% %   di MCNPX, effettuato in un altro studio.
% %  \begin{columns}
% %   \begin{column}{0.5\textwidth}
% %    %%%%%%%%%%%%%%%%%%%%%%%%%%%%%%%%%%%%%%%%%%%%%%%%%%%%%%%%%%%%%%
% % %   \begin{equation*}
% % %    k_{eff}=\frac{popolazione neutronica del}{×}
% % %   \end{equation*}
% % 
% % 
% %   \end{column}
% % 
% %   \begin{column}{0.4\textwidth}
% %   \begin{center}
% %       \begin{table}
% %  \begin{tabular}{lcc}
% % \hline
% % Tempo  & Donjon  &  MCNPX \\
% % $[$giorni$]$ & $k_{eff}$ [-] & $k_{eff}$ [-]  \\
% % \hline
% % 0   & 1.07935 & 1.0525 \\
% % 365 & 1.0520  & 1.0229  \\
% % 730 & 1.0272  & 0.9964  \\
% % 1095 & 1.0001 & 0.9703  \\
% % \hline
% % \end{tabular}
% % \end{table}
% % %%%%%%%%%%%%%%%%%%%%%%%%%%%%%%%%%%%%%%%%%%%%%%%%%%%%%%%%%%%%%%
% %   \end{center}
% % 
% %   \end{column}
% %   
% %  \end{columns}
% % 
% % \end{frame}
% % % % % % % % % % %%%%%%%%%%%%%%%%%%%%%%%%%%%%%%%%%%%%%%%%%%%%%%%%%%%%%%%%%%%%%%%%%%%%%%%%%%%%%%%%
% % % % % % % % % % \begin{frame}
% % % % % % % % % % \frametitle{Variazioni di $k_{eff}$ in funzione della temperatura $T_f$}
% % % % % % % % % %   \begin{center}
% % % % % % % % % %    \includegraphics[scale=0.7]{../../thesis/graph/keff_temp.eps}
% % % % % % % % % %   \end{center}
% % % % % % % % % % \end{frame}
% % % % % % % % % % %%%%%%%%%%%%%%%%%%%%%%%%%%%%%%%%%%%%%%%%%%%%%%%%%%%%%%%%%%%%%%
% % % % % % % % % % %%%%%%%%%%%%%%%%%%%%%%%%%%%%%%%%%%%%%%%%%%%%%%%%%%%%%%%%%%%%%%
\begin{frame}
\frametitle{Variazioni di $k_{eff}$ in funzione dei tempi di bruciamento $t_{burnup}$}
\vspace{0.5cm}
  \begin{center}
   \includegraphics[scale=0.7]{../../thesis/graph/keff_in_out_side0-eps-converted-to}
  
  \begin{itemize}
 \item \centering Confronto tra barre inserite e estratte per temperatura 
 del combustibile $T_f$= 550 $^\circ$C e $T_f$= 1000 $^\circ$C 
\end{itemize}
\end{center}
\end{frame}
%%%%%%%%%%%%%%%%%%%%%%%%%%%%%%%%%%%%%%%%%%%%%%%%%%%%%%%%%%%%%%
%%%%%%%%%%%%%%%%%%%%%%%%%%%%%%%%%%%%%%%%%%%%%%%%%%%%%%%%%%%%%%%%%%%%%%%%%%%%%%%%
\begin{frame}
\frametitle{Distribuzione di flusso neutronico a $t_{burnup}$= 3 anni}
    \vspace{0.5cm}
 \begin{columns}
  \begin{column}{0.48\textwidth}
  \begin{center}
  \underline{Barre estratte}
  
   \includegraphics[scale=0.2]{../../thesis/figure/result_tred/morenew/flux-trashold-exstract1}
   
  \end{center}

  
  \end{column}

  \begin{column}{0.48\textwidth}
  
  \begin{center}
  \underline{Barre inserite}
  
   \includegraphics[scale=0.2]{../../thesis/figure/result_tred/morenew/flux-trashold-insert1}
   
  \end{center}
  
  \end{column}
  
 \end{columns}
% 
% 
\end{frame}
%%%%%%%%%%%%%%%%%%%%%%%%%%%%%%%%%%%%%%%%%%%%%%%%%%%%%%%%%%%%%%
%%%%%%%%%%%%%%%%%%%%%%%%%%%%%%%%%%%%%%%%%%%%%%%%%%%%%%%%%%%%%%%%%%%%%%%%%%%%%%%%
\begin{frame}
\frametitle{Distribuzione di potenza termica a $t_{burnup}$= 3 anni}
  \vspace{0.5cm}
 \begin{columns}
  \begin{column}{0.48\textwidth}
  
  \begin{center}
  \underline{Barre estratte}
  
   \includegraphics[scale=0.202]{../../thesis/figure/result_tred/morenew/pw-trashold-exstract1}
  \end{center}

  
  \end{column}

  \begin{column}{0.48\textwidth}
  
  \begin{center}
  \underline{Barre inserite}
  
   \includegraphics[scale=0.2]{../../thesis/figure/result_tred/morenew/pw-trashold-insert1}
  \end{center}
  
  \end{column}
  
 \end{columns}
% 
% 
\end{frame}
%%%%%%%%%%%%%%%%%%%%%%%%%%%%%%%%%%%%%%%%%%%%%%%%%%%%%%%%%%%%%%
% % % % %%%%%%%%%%%%%%%%%%%%%%%%%%%%%%%%%%%%%%%%%%%%%%%%%%%%%%%%%%%%%%%%%%%%%%%%%%%%%%%%
% % % % \begin{frame}
% % % % \frametitle{Confronto barre inserite e barre estratte lungo l'asse $y$}
% % % %   \vspace{0.5cm}
% % % %   
% % % %  \begin{columns}
% % % %   \begin{column}{0.48\textwidth}
% % % %   
% % % %   \begin{center}
% % % %   Distribuzione di flusso
% % % %    \includegraphics[scale=0.43]{../../thesis/graph/fluxy-eps-converted-to}
% % % %   \end{center}
% % % % 
% % % %   
% % % %   \end{column}
% % % % 
% % % %   \begin{column}{0.48\textwidth}
% % % %   
% % % %   \begin{center}
% % % %   Distribuzione di potenza termica
% % % %    \includegraphics[scale=0.43]{../../thesis/graph/powery}
% % % %   \end{center}
% % % %   
% % % %   \end{column}
% % % %   
% % % %  \end{columns}
% % % %  \vspace{0.4cm}
% % % % %         \hspace{0.5cm}
% % % %  \begin{center}
% % % %  
% % % %  \begin{block}{Relazione tra flusso neutronico e potenza termica}
% % % % %  \hspace{1.5cm}
% % % %   \begin{equation*}
% % % %    P_{th}=E_{fissione} \cdot V_{nocciolo} \cdot \int_E \Sigma_F(E) \Phi(E)\, dE
% % % %   \end{equation*}
% % % % 
% % % %  \end{block}
% % % %   \end{center}
% % % % % 
% % % % \end{frame}
% % % % %%%%%%%%%%%%%%%%%%%%%%%%%%%%%%%%%%%%%%%%%%%%%%%%%%%%%%%%%%%%%%%%%%%%%%%%%%%%5
%%%%%%%%%%%%%%%%%%%%%%%%%%%%%%%%%%%%%%%%%%%%%%%%%%%%%%%%%%%%%%%%%%%%%%%%%%%%%%%%
\begin{frame}
\frametitle{Confronto barre inserite e barre estratte lungo un profilo radiale}
  \vspace{0.5cm}
  
 \begin{columns}
  \begin{column}{0.48\textwidth}
  
  \begin{center}
  Distribuzione di flusso
     \includegraphics[scale=0.43]{../../thesis/graph/fluxy-eps-converted-to}
%    \includegraphics[height=0.38\textheight]{../flusso}
  \end{center}

  
  \end{column}
 \hspace{-1cm}
  \begin{column}{0.48\textwidth}
  
  \begin{center}
  Distribuzione di potenza termica
   \includegraphics[scale=0.43]{../potenza}
  \end{center}
  
  \end{column}
  
 \end{columns}
 \vspace{0.4cm}
%         \hspace{0.5cm}
 \begin{center}
 
 \begin{block}{Relazione tra flusso neutronico e potenza termica}
%  \hspace{1.5cm}
  \begin{equation*}
   P_{th}=E_{fissione} \cdot V_{nocciolo} \cdot \int_E \Sigma_F(E) \Phi(E)\, dE
  \end{equation*}

 \end{block}
  \end{center}
% 
\end{frame}
%%%%%%%%%%%%%%%%%%%%%%%%%%%%%%%%%%%%%%%%%%%%%%%%%%%%%%%%%%%%%%%%%%%%%%%%%%%%%%5
\begin{frame}
\frametitle{Evoluzione del campo di temperatura}
\vspace{0.3cm} 
 \begin{columns}
  \begin{column}{0.3\textwidth}
  
  \begin{center}
  \underline{$t$= 100 secondi}
  \vspace{0.1cm}
  
   \includegraphics[scale=0.2]{../../thesis/figure/result_tred/new/tred2-temp}
   
  \end{center}

  
  \end{column}

  \begin{column}{0.3\textwidth}
  
  \begin{center}
   \underline{$t$= 250 secondi}
   \vspace{0.1cm}
   
   \includegraphics[scale=0.2]{../../thesis/figure/result_tred/new/tred3-temp}
  
  \end{center}
  
  \end{column}
  
   \begin{column}{0.3\textwidth}
  
  \begin{center}
   \underline{$t$= 400 secondi}
   \vspace{0.1cm}
   
   \includegraphics[scale=0.2]{../../thesis/figure/result_tred/new/tred1-temp}
  
  \end{center}
  
  \end{column}
 
 \end{columns}
 \vspace{0.3cm}
 
 \begin{center}
    Velocità del piombo liquido considerata costante \\
    \textbf{$\boldsymbol{\upsilon}$ = 1.28 m/s}
 \end{center}


% 
% 
\end{frame}
%%%%%%%%%%%%%%%%%%%%%%%%%%%%%%%%%%%%%%%%%%%%%%%%%%%%%%%%%%%%%%
%%%%%%%%%%%%%%%%%%%%%%%%%%%%%%%%%%%%%%%%%%%%%%%%%%%%%%%%%%%%%%
% % % % % % % % % % % % % % % % \begin{frame}
% % % % % % % % % % % % % % % % \frametitle{Distribuzione di temperatura lungo l'asse $z$}
% % % % % % % % % % % % % % % %   
% % % % % % % % % % % % % % % %  \begin{columns}
% % % % % % % % % % % % % % % %   \begin{column}{0.4\textwidth}
% % % % % % % % % % % % % % % %   
% % % % % % % % % % % % % % % %   \begin{center}
% % % % % % % % % % % % % % % %    \includegraphics[scale=0.2]{../../thesis/figure/result_tred/morenew/tempassurda}
% % % % % % % % % % % % % % % %   \end{center}
% % % % % % % % % % % % % % % % 
% % % % % % % % % % % % % % % %   
% % % % % % % % % % % % % % % %   \end{column}
% % % % % % % % % % % % % % % % 
% % % % % % % % % % % % % % % %   \begin{column}{0.4\textwidth}
% % % % % % % % % % % % % % % %   
% % % % % % % % % % % % % % % %   \begin{center}
% % % % % % % % % % % % % % % %   Profilo di temperatura
% % % % % % % % % % % % % % % %    \includegraphics[scale=0.4]{../../thesis/graph/tempz}
% % % % % % % % % % % % % % % %   \end{center}
% % % % % % % % % % % % % % % %   
% % % % % % % % % % % % % % % %   \end{column}
% % % % % % % % % % % % % % % %   
% % % % % % % % % % % % % % % %  \end{columns}
% % % % % % % % % % % % % % % % % 
% % % % % % % % % % % % % % % % % 
% % % % % % % % % % % % % % % % \end{frame}
% % % % % % %%%%%%%%%%%%%%%%%%%%%%%%%%%
% % % % % % \section{Conclusioni}
% % % % % % %%%%%%%%%%%%%%%%%%%%%%%%%%%%%%%%%%%%%%%%%%%%%%%%%%%%%%%%%%%%%%%%%%%%%%%%%%%%%%%%%%%%%%%%%%%%%%%%%%%%%%%%%%
% % % % % % \begin{frame}
% % % % % % \frametitle{Conclusioni}
% % % % % % \begin{block}{Aspetti positivi}
% % % % % % \begin{itemize}
% % % % % % \item L'approccio multifisico consente di effettuare una simulazione complessiva della fisica del reattore
% % % % % % \item L'evoluzione isotopica dovuta al bruciamento è in linea con altri studi
% % % % % % \item L'accoppiamento neutronico-termoidraulico che è stato implementato apre a prospettive future per l'indagine del compormento multifisico
% % % % % % di sistemi nucleari
% % % % % % %\item La soluzione è ottenuta dopo numerose iterazione attraverso il codice di neutronica e termoidraulica con correzioni di temperatura e densità
% % % % % % %\item I risultati ottenuti aprono a prosfettivi di estendere un simile approccio ad altri modelli
% % % % % % %\item L'intera procedura è in grado di garantire flessibilità all'utente
% % % % % % \end{itemize}
% % % % % % \end{block}
% % % % % % \begin{block}{Aspetti negativi}
% % % % % % \begin{itemize}
% % % % % % \item $k_{eff}$ troppo elevato in riferimento ad altri studi del settore
% % % % % % \item Possibilità di interpolare solo su due variabili e non su tre 
% % % % % % \item Modello di termoidraulica troppo semplificativo per un'analisi approffondita
% % % % % % \end{itemize}
% % % % % % \end{block}
% % % % % % 
% % % % % % \end{frame}
% % % % % % 
% % % % % % % \section{Sviluppi Futuri}
% % % % % % \begin{frame}
% % % % % % \frametitle{Sviluppi Futuri}
% % % % % % \begin{itemize}
% % % % % % \item Aggiornamento dei dati utilizzati nella modellazione con i progetti più recenti inerenti al progetto ALFRED
% % % % % % \item Approfondimento sui sistemi di sicurezza del reattore nel codice di neutronica
% % % % % % \item Raffinazione del modello termoidraulico attraverso l'introduzione di modelli più sosfisticati e l'implementazione di codici per la turbolenza dei metalli liquidi
% % % % % % \item Estendere un simile approccio multifisico ad altri reattori
% % % % % % \end{itemize}
% % % % % % \end{frame}
% % % % % % %%%%%%%%%%%%%%%%%%%%%%%%%%%%%%%%%%%%%%%%%%%%%%%%%%%%%%%%%%%%%%%%%%%%%%%%%%%%%%%%%%%%%%%%%%%%%%%%%%%%%%%%%%
